\documentclass[11pt]{article}
\usepackage[margin=1in]{geometry}
\usepackage{amsmath, amssymb}
\usepackage{graphicx}
\usepackage{booktabs}
\usepackage{enumitem}
\usepackage{fancyhdr}
\usepackage{tcolorbox}

\pagestyle{fancy}
\fancyhf{}
\rhead{Sample Quiz 2 - Solutions}
\lhead{PHYSICS 1101.100}
\rfoot{Page \thepage}

\tcbuselibrary{theorems, skins}

\newtcolorbox{keybox}[1][]{
  colback=yellow!10,
  colframe=orange!75!black,
  title=#1,
  fonttitle=\bfseries
}

\newtcolorbox{solutionbox}{
  colback=blue!5,
  colframe=blue!50!black,
  title=Solution,
  fonttitle=\bfseries
}

\begin{document}

\begin{center}
{\LARGE \textbf{PHYSICS 1101.100 -- Sample Quiz 2}}\\[0.5em]
{\large Additional Practice Problems: Complete Solutions}
\end{center}

\vspace{1em}

%=============================================
\section*{Practice Problem A: Two-Spring System Variation}
%=============================================

\begin{keybox}[Problem Statement]
A diving board of length $L = 3.0$ m and mass $M = 25$ kg is supported by two springs. Spring 1 (with constant $k_1 = 8000$ N/m) is located at the left end, and Spring 2 (with constant $k_2 = 12000$ N/m) is located at a distance of 2.0 m from the left end. A diver of mass $m = 70$ kg stands at the right end of the board. Find the compression of each spring.
\end{keybox}

\begin{solutionbox}
\textbf{Step 1: Identify all forces and their positions}

Let $\delta_1$ and $\delta_2$ be the compressions of springs 1 and 2.

\begin{itemize}
    \item Spring 1 force: $F_1 = k_1\delta_1$ at $x = 0$
    \item Spring 2 force: $F_2 = k_2\delta_2$ at $x = 2.0$ m
    \item Board weight: $W_{\text{board}} = Mg = (25)(9.8) = 245$ N at $x = 1.5$ m
    \item Diver weight: $W_{\text{diver}} = mg = (70)(9.8) = 686$ N at $x = 3.0$ m
\end{itemize}

\textbf{Step 2: Apply Force Equilibrium ($\sum F_y = 0$)}
\begin{equation}
k_1\delta_1 + k_2\delta_2 = 245 + 686 = 931 \text{ N}
\end{equation}

\textbf{Step 3: Apply Torque Equilibrium ($\sum \tau = 0$ about left end)}
\begin{align}
k_2\delta_2(2.0) &= W_{\text{board}}(1.5) + W_{\text{diver}}(3.0)\\
24000\delta_2 &= 367.5 + 2058 = 2425.5\\
\delta_2 &= \boxed{0.101 \text{ m} = 10.1 \text{ cm}}
\end{align}

\textbf{Step 4: Solve for $\delta_1$}
\begin{align}
8000\delta_1 + 12000(0.101) &= 931\\
8000\delta_1 &= 931 - 1212 = -281\\
\delta_1 &= \boxed{-0.035 \text{ m}}
\end{align}

The negative sign indicates Spring 1 is in \textbf{tension} (stretched by 3.5 cm).
\end{solutionbox}

\vspace{2em}

%=============================================
\section*{Practice Problem B: Variable-Force Stopping}
%=============================================

\begin{keybox}[Problem Statement]
A 1500 kg car traveling at 25 m/s needs to stop. The braking system applies a force that varies with position according to $F(x) = F_0 + bx$, where $F_0 = 5000$ N and $b = 200$ N/m. Using energy methods, find the stopping distance and compare it to constant-force braking with force $F_0$.
\end{keybox}

\begin{solutionbox}
\textbf{Step 1: Calculate Initial Kinetic Energy}
\begin{equation}
KE_i = \frac{1}{2}mv_0^2 = \frac{1}{2}(1500)(25)^2 = 468,750 \text{ J}
\end{equation}

\textbf{Step 2: Calculate Work by Variable Force}
\begin{equation}
W = -\int_0^{d}(F_0 + bx)\,dx = -\left(F_0 d + \frac{bd^2}{2}\right)
\end{equation}

\textbf{Step 3: Apply Work-Energy Theorem}
\begin{align}
F_0 d + \frac{bd^2}{2} &= 468,750\\
5000d + 100d^2 &= 468,750\\
d^2 + 50d - 4687.5 &= 0
\end{align}

\textbf{Step 4: Solve Quadratic}
\begin{equation}
d = \frac{-50 + \sqrt{2500 + 18750}}{2} = \frac{-50 + 145.8}{2} = \boxed{47.9 \text{ m}}
\end{equation}

\textbf{Step 5: Constant-Force Comparison}
\begin{equation}
d_{\text{constant}} = \frac{KE_i}{F_0} = \frac{468,750}{5000} = \boxed{93.75 \text{ m}}
\end{equation}
\end{solutionbox}

\vspace{1em}

\begin{keybox}[Key Result]
\begin{center}
\begin{tabular}{lc}
\toprule
\textbf{Braking Type} & \textbf{Stopping Distance}\\
\midrule
Variable Force ($F_0 + bx$) & 47.9 m\\
Constant Force ($F_0$) & 93.75 m\\
\bottomrule
\end{tabular}
\end{center}
\textbf{Conclusion:} Variable-force braking reduces stopping distance by \textbf{49\%}!
\end{keybox}

%=============================================
\section*{Mastery Check: Practice Questions}
%=============================================

Test your understanding with these variations:

\begin{enumerate}[label=\textbf{Q\arabic*.}]
    \item What if the braking force were $F(x) = F_0 - bx$ (decreasing with distance)?
    \item In Problem A, where should Spring 2 be placed to make both compressions equal?
    \item If the car's mass doubled, how would the stopping distance change for each case?
\end{enumerate}

\end{document}
```

---

## �� Mastery Checklist

Before your test, make sure you can:

- [ ] Set up torque equilibrium equations with correct moment arms
- [ ] Integrate variable forces to find work
- [ ] Apply the work-energy theorem correctly
- [ ] Solve quadratic equations from physics problems
- [ ] Interpret negative values physically (tension vs. compression)
- [ ] Compare results to build physical intuition

Would you like me to create additional practice variations or explain any step in more detail?
%=============================================
\section*{Mastery Check: Detailed Solutions}
%=============================================

\begin{keybox}[Solutions to Practice Questions]
\textbf{Q1. Decreasing Force $F(x) = F_0 - bx$}

The work done becomes:
\begin{equation}
W = -\int_0^{d}(F_0 - bx)\,dx = -\left(F_0 d - \frac{bd^2}{2}\right)
\end{equation}

Setting $|W| = KE_i$:
\begin{equation}
5000d - 100d^2 = 468,750
\end{equation}
\begin{equation}
100d^2 - 5000d + 468,750 = 0 \quad \Rightarrow \quad d^2 - 50d + 4687.5 = 0
\end{equation}

Using the quadratic formula:
\begin{equation}
d = \frac{50 \pm \sqrt{2500 - 18750}}{2} = \frac{50 \pm \sqrt{-16250}}{2}
\end{equation}

\textbf{Result:} No real solution! The decreasing force becomes too weak to stop the car. The force would reach zero at $x = F_0/b = 25$ m, but the car still has kinetic energy remaining.

\vspace{1em}

\textbf{Q2. Equal Compression Condition}

For $\delta_1 = \delta_2 = \delta$, we need:
\begin{equation}
\frac{F_1}{k_1} = \frac{F_2}{k_2} \quad \Rightarrow \quad \frac{F_1}{F_2} = \frac{k_1}{k_2} = \frac{8000}{12000} = \frac{2}{3}
\end{equation}

Let Spring 2 be at distance $a$ from the left end. Using torque about Spring 1:
\begin{equation}
F_2 \cdot a = W_{\text{board}} \cdot \frac{L}{2} + W_{\text{diver}} \cdot L
\end{equation}
\begin{equation}
F_2 \cdot a = (245)(1.5) + (686)(3.0) = 367.5 + 2058 = 2425.5 \text{ N·m}
\end{equation}

From force balance: $F_1 + F_2 = 931$ N, and with $F_1 = \frac{2}{3}F_2$:
\begin{equation}
\frac{2}{3}F_2 + F_2 = 931 \quad \Rightarrow \quad F_2 = 558.6 \text{ N}
\end{equation}

Therefore:
\begin{equation}
a = \frac{2425.5}{558.6} = \boxed{4.34 \text{ m}}
\end{equation}

\textbf{Note:} This is beyond the board length (3.0 m), so equal compression is \textbf{not physically achievable} with this configuration!

\vspace{1em}

\textbf{Q3. Doubled Mass Effect}

If $m \to 2m = 3000$ kg:
\begin{itemize}
    \item $KE_i \to 2 \times 468,750 = 937,500$ J
\end{itemize}

\textit{Constant force case:}
\begin{equation}
d_{\text{constant}} = \frac{937,500}{5000} = \boxed{187.5 \text{ m}} \quad \text{(exactly doubled)}
\end{equation}

\textit{Variable force case:}
\begin{equation}
100d^2 + 5000d = 937,500 \quad \Rightarrow \quad d^2 + 50d - 9375 = 0
\end{equation}
\begin{equation}
d = \frac{-50 + \sqrt{2500 + 37500}}{2} = \frac{-50 + 200}{2} = \boxed{75 \text{ m}}
\end{equation}

\textbf{Key Insight:} The variable force case increases by factor of $75/47.9 = 1.57$, \textit{less than doubled}, because the increasing force provides more braking at larger distances.
\end{keybox}

%=============================================
\section*{Conceptual Understanding: Key Physics Principles}
%=============================================

\subsection*{Principle 1: Static Equilibrium}

\begin{conceptbox}[Equilibrium Conditions]
For any rigid body in static equilibrium:
\begin{enumerate}
    \item \textbf{Force Balance:} $\sum \vec{F} = 0$
    \item \textbf{Torque Balance:} $\sum \vec{\tau} = 0$ (about \textit{any} point)
\end{enumerate}

\textbf{Strategy for choosing pivot point:} Select the pivot at the location of an unknown force to eliminate it from the torque equation!
\end{conceptbox}

\begin{center}
\begin{tikzpicture}[scale=1.2]
    % Board
    \draw[thick, fill=brown!30] (0,0) rectangle (6,0.2);
    
    % Springs
    \draw[thick, decorate, decoration={coil, segment length=3pt, amplitude=4pt}] (0.5,-0.8) -- (0.5,0);
    \draw[thick, decorate, decoration={coil, segment length=3pt, amplitude=4pt}] (4,-0.8) -- (4,0);
    
    % Ground
    \draw[thick] (-0.5,-0.8) -- (6.5,-0.8);
    \fill[pattern=north east lines] (-0.5,-1) rectangle (6.5,-0.8);
    
    % Forces
    \draw[->, thick, blue] (0.5,0.2) -- (0.5,1.2) node[above] {$F_1$};
    \draw[->, thick, blue] (4,0.2) -- (4,1.5) node[above] {$F_2$};
    \draw[->, thick, red] (3,-0.1) -- (3,-1.1) node[below] {$Mg$};
    \draw[->, thick, red] (5.8,-0.1) -- (5.8,-1.5) node[below] {$mg$};
    
    % Labels
    \node at (0.5,-1.1) {Spring 1};
    \node at (4,-1.1) {Spring 2};
    \draw[<->] (0,-1.5) -- (4,-1.5) node[midway, below] {$a$};
    \draw[<->] (0,-2) -- (6,-2) node[midway, below] {$L$};
    
    % Pivot indicator
    \draw[fill=green!50] (0.5,0) circle (0.1);
    \node[green!50!black] at (0.5,0.5) {Pivot};
\end{tikzpicture}
\end{center}

\subsection*{Principle 2: Work-Energy Theorem}

\begin{conceptbox}[Work-Energy Theorem]
\begin{equation}
W_{\text{net}} = \Delta KE = KE_f - KE_i
\end{equation}

For a variable force $F(x)$:
\begin{equation}
W = \int_{x_i}^{x_f} F(x)\, dx
\end{equation}

\textbf{Sign Convention:}
\begin{itemize}
    \item Force in direction of motion: $W > 0$ (adds energy)
    \item Force opposite to motion: $W < 0$ (removes energy)
\end{itemize}
\end{conceptbox}

\subsection*{Principle 3: Spring Mechanics}

\begin{conceptbox}[Hooke's Law and Spring Energy]
\textbf{Hooke's Law:}
\begin{equation}
F_{\text{spring}} = -kx \quad \text{(restoring force)}
\end{equation}
\begin{equation}
F_{\text{applied}} = kx \quad \text{(force needed to compress/stretch)}
\end{equation}

\textbf{Elastic Potential Energy:}
\begin{equation}
U_{\text{spring}} = \frac{1}{2}kx^2
\end{equation}

\textbf{Work to compress a spring:}
\begin{equation}
W = \int_0^{\delta} kx\, dx = \frac{1}{2}k\delta^2
\end{equation}
\end{conceptbox}

%=============================================
\section*{Problem-Solving Framework}
%=============================================

\begin{keybox}[Step-by-Step Strategy for Equilibrium Problems]
\begin{enumerate}
    \item \textbf{Draw a Free Body Diagram}
    \begin{itemize}
        \item Identify all forces (gravity, springs, applied forces)
        \item Mark distances and positions clearly
    \end{itemize}
    
    \item \textbf{Choose a Coordinate System}
    \begin{itemize}
        \item Define positive directions for forces
        \item Choose a convenient pivot point for torques
    \end{itemize}
    
    \item \textbf{Write Equilibrium Equations}
    \begin{itemize}
        \item $\sum F_y = 0$ (vertical force balance)
        \item $\sum \tau = 0$ (torque balance about chosen pivot)
    \end{itemize}
    
    \item \textbf{Solve the System}
    \begin{itemize}
        \item Use substitution or matrix methods
        \item Check units at each step
    \end{itemize}
    
    \item \textbf{Verify and Interpret}
    \begin{itemize}
        \item Do the answers make physical sense?
        \item Check limiting cases
    \end{itemize}
\end{enumerate}
\end{keybox}

\begin{keybox}[Step-by-Step Strategy for Work-Energy Problems]
\begin{enumerate}
    \item \textbf{Identify Initial and Final States}
    \begin{itemize}
        \item What is the initial kinetic energy?
        \item What is the final kinetic energy? (often zero for stopping problems)
    \end{itemize}
    
    \item \textbf{Identify All Forces Doing Work}
    \begin{itemize}
        \item Is the force constant or variable?
        \item What is the direction relative to motion?
    \end{itemize}
    
    \item \textbf{Calculate Work}
    \begin{itemize}
        \item Constant force: $W = Fd\cos\theta$
        \item Variable force: $W = \int F(x)\, dx$
    \end{itemize}
    
    \item \textbf{Apply Work-Energy Theorem}
    \begin{itemize}
        \item Set $W_{\text{net}} = \Delta KE$
        \item Solve for unknown (distance, velocity, etc.)
    \end{itemize}
\end{enumerate}
\end{keybox}

%=============================================
\section*{Common Mistakes to Avoid}
%=============================================

\begin{warningbox}[Critical Errors on Exams]
\begin{enumerate}
    \item \textbf{Torque Sign Errors}
    \begin{itemize}
        \item Always use consistent sign convention (CCW = positive)
        \item Double-check which forces create CW vs CCW rotation
    \end{itemize}
    
    \item \textbf{Forgetting the Board's Weight}
    \begin{itemize}
        \item The board's weight acts at its center of mass ($L/2$)
        \item Don't ignore it even when a person is also on the board
    \end{itemize}
    
    \item \textbf{Integration Limits}
    \begin{itemize}
        \item For stopping problems: integrate from 0 to $d$ (stopping distance)
        \item Watch the sign of the integrand
    \end{itemize}
    
    \item \textbf{Quadratic Formula Errors}
    \begin{itemize}
        \item Always check: is $d^2 + 50d - 4687.5 = 0$ or $d^2 + 50d = 4687.5$?
        \item Physical solutions must be positive
    \end{itemize}
    
    \item \textbf{Unit Inconsistencies}
    \begin{itemize}
        \item Ensure all quantities are in SI units before calculating
        \item N/m, kg, m/s, J must all be consistent
    \end{itemize}
\end{enumerate}
\end{warningbox}

%=============================================
\section*{Quick Reference Formulas}
%=============================================

\begin{center}
\begin{tabular}{|l|c|l|}
\hline
\textbf{Concept} & \textbf{Formula} & \textbf{Units} \\
\hline
\hline
Kinetic Energy & $KE = \frac{1}{2}mv^2$ & J (Joules) \\
\hline
Spring Force & $F = kx$ & N (Newtons) \\
\hline
Spring Potential Energy & $U = \frac{1}{2}kx^2$ & J \\
\hline
Work (constant force) & $W = Fd\cos\theta$ & J \\
\hline
Work (variable force) & $W = \int F(x)\,dx$ & J \\
\hline
Torque & $\tau = rF\sin\theta$ & N·m \\
\hline
Weight & $W = mg$ & N \\
\hline
Quadratic Formula & $x = \frac{-b \pm \sqrt{b^2-4ac}}{2a}$ & -- \\
\hline
\end{tabular}
\end{center}

%=============================================
\section*{Self-Assessment: Are You Ready?}
%=============================================

\begin{conceptbox}[Pre-Test Checklist]
Rate yourself 1-5 on each skill (5 = completely confident):

\begin{tabular}{|p{10cm}|c|}
\hline
\textbf{Skill} & \textbf{Rating} \\
\hline
I can draw accurate free body diagrams & \quad /5 \\
\hline
I can set up torque equations with correct signs & \quad /5 \\
\hline
I can solve systems of two equations & \quad /5 \\
\hline
I can integrate polynomial functions & \quad /5 \\
\hline
I can apply the work-energy theorem & \quad /5 \\
\hline
I can solve quadratic equations & \quad /5 \\
\hline
I can interpret physical meaning of results & \quad /5 \\
\hline
\end{tabular}

\vspace{1em}
\textbf{If any rating is below 4, review that section before the test!}
\end{conceptbox}

%=============================================
\section*{Additional Practice: Challenge Problems}
%=============================================

\begin{solutionbox}
\textbf{Challenge 1: Three-Spring System}

A 4.0 m beam (mass 30 kg) is supported by three springs:
\begin{itemize}
    \item Spring A ($k_A = 6000$ N/m) at x = 0
    \item Spring B ($k_B = 10000$ N/m) at x = 2.0 m
    \item Spring C ($k_C = 8000$ N/m) at x = 4.0 m
\end{itemize}

A 50 kg mass is placed at x = 3.0 m. Find all three spring compressions.

\textit{Hint: You need three equations - two torque equations (about different points) and one force equation.}
\end{solutionbox}

\begin{solutionbox}
\textbf{Challenge 2: Velocity-Dependent Braking}

A 2000 kg car traveling at 30 m/s experiences a braking force $F(v) = cv^2$ where $c = 50$ N·s²/m². 

Find the stopping distance using $F = ma = m\frac{dv}{dt} = mv\frac{dv}{dx}$.

\textit{Hint: Separate variables and integrate with respect to velocity.}
\end{solutionbox}

\vspace{2em}
\begin{center}
\fbox{\parbox{0.8\textwidth}{
\centering
\Large\textbf{Good Luck on Your Quiz!}\\[0.5em]
\normalsize Remember: Draw diagrams, check units, and verify your answers make physical sense.
}}
\end{center}

\end{document}
